\chapter*{Abstract}
\addcontentsline{toc}{chapter}{Abstract}

The cybersecurity landscape has evolved dramatically from simple proof-of-concept viruses to sophisticated criminal operations targeting financial gain and critical infrastructure. Modern malware employs advanced evasion techniques, including obfuscation, packing, steganography, and binary padding, to bypass traditional security defenses. This research investigates contemporary malware techniques from an offensive security perspective to understand how attackers build, deliver, and maintain malicious software.

Through controlled laboratory experiments, this study examines five key areas: steganographic delivery through malicious PDFs, software packing using UPX, binary padding to evade size-based scanning limits, persistence mechanisms via Windows services, and the effectiveness of signature-based, heuristic-based, and behaviour-based detection methods. The research demonstrates that low-complexity techniques can achieve high-impact results, with traditional signature-based detection proving inadequate against structurally modified malware variants.

Key findings reveal that steganography provides a simple yet powerful delivery mechanism, packing easily defeats signature-based antivirus solutions, and binary padding exploits performance-based scanning weaknesses in commercial security tools. The study concludes that defensive strategies must prioritize behavioural detection, enhance email security and user awareness, implement monitoring of persistence mechanisms, and adopt multi-layered security approaches. These insights provide practical guidance for security analysts, incident responders, SOC teams, and cybersecurity researchers working to strengthen organizational defenses against evolving malware threats.

\textbf{Keywords:} Malware analysis, obfuscation techniques, evasion methods, steganography, software packing, cybersecurity, offensive security, threat detection
